\subsection*{The actual compact quartic action and its nonzero CM component}
\begin{statement}{E.109}{Proved}{The compact class-two arithmetic action, its exact affine comparison, and its entire CM projection}
Retain the plus quartic $C$, its affine open $U$, the field
$F=\mathbb Q(s)$ with $s=\zeta _8$, and the actual class-two
section $S_5$ of E.105.  All groups in this statement are the
geometric pro-five quotients of those arithmetic fundamental groups,
with continuous $\mathbb Z_5$-module cohomology.  Put
\[
 R=\mathbb Z_5,\quad A=T_5J_{C,D},\quad H=T_5J_C,
 \quad B=\Lambda^2 A,\quad Z=\Lambda^2H/Rq,
 \quad L=(\Lambda^3H)(-1),
 \tag{E109.1}
\]
where the polarization map $R(1)\to\Lambda^2H$ has image $Rq$.
The ranks of $A,B,H,Z,L$ are $17,136,6,14,20$.
Write
\[
 \mathcal D=\operatorname{Hom}_R(H,Z)/\operatorname{ad}H,
 \qquad \operatorname{ad}_h(x)=h\wedge x\pmod{Rq}.
 \tag{E109.2}
\]
It is free of rank $78$.  The symplectic Johnson map is the actual
$G_{\mathbb Q}$-equivariant map
\[
 f:L\longrightarrow\operatorname{Hom}_R(H,Z),\qquad
 f(a\wedge b\wedge c)(x)=
 \langle a,x\rangle b\wedge c+
 \langle b,x\rangle c\wedge a+
 \langle c,x\rangle a\wedge b\pmod{Rq}.
 \tag{E109.3}
\]
The twist in $L$ cancels the value in $R(1)$ of each pairing.
Since $f(q\wedge h)=-\operatorname{ad}_h$, it induces an
injection $\bar f:L/(q\wedge H)\to\mathcal D$.

Choose the square root $i=s^2$ in the CM action; the two embeddings
of $\mathbb Z[i]$ into $R$ define $G_F$-stable summands
$H=H^+\oplus H^-$ of rank three.  The module
\[
 M=(\Lambda^3H^+\oplus\Lambda^3H^-)(-1)
       \ \subset L/(q\wedge H)
 \tag{E109.4}
\]
has rank two and is stable under $G_{\mathbb Q}$, which can
interchange the two summands.  There is a continuous
$G_{\mathbb Q}$-equivariant left inverse
\[
 \Pi:\mathcal D\longrightarrow M,
 \qquad \Pi\bar f|_M=1_M.
 \tag{E109.5}
\]
Formula (E109.10) below gives it on every coefficient.

In BCH coordinates let $c^C$ denote the normalized central
off-diagonal cocycle of the actual compact arithmetic action,
and let $[c^C]_{\rm out}\in H^1_{\rm cts}(G_{\mathbb Q},\mathcal D)$
be its image.  There is a specified class
$\beta_{\mathbb Q}=2\Pi_*[c^C]_{\rm out}$ in
$H^1_{\rm cts}(G_{\mathbb Q},M)$ such that
\[
 2[c^C]_{\rm out}=\bar f_*\beta_{\mathbb Q},\qquad
 \beta_{\mathbb Q}|_{G_F}=\pi_{M*}\mu(C,b),\qquad
 b=[s:0:1],\qquad
 \beta_{\mathbb Q}\otimes_R\mathbb Q_5\ne0.
 \tag{E109.6}
\]
Here $\mu(C,b)$ is the integral five-adic Abel--Jacobi class
of $i_b(C)-[-1]i_b(C)$, and $\pi_M$ is the displayed pure-CM
projector.  In particular (E109.6) determines the entire outer
action class through this rank-two module.  The explicit
finite-field certificate uses $p=19$:
\[
 E':Y^2=X^3+X,\qquad
 \delta_{19}=(17,16),\quad \operatorname{ord}(\delta_{19})=10,
 \quad 4\delta_{19}=(9,4),\quad
 \operatorname{ord}(4\delta_{19})=5.
 \tag{E109.7}
\]

The source of this class in the actual affine action is also
explicit.  Let $p:A\to H$ be filling the cusps, let
$p_B=\Lambda^2p\pmod{Rq}$, and let $j:H\to A$ be the actual
equivariant section $j_*/4$ from E.103.  The central cocycle
$c^A\in Z^1_{\rm cts}(G_{\mathbb Q},\operatorname{Hom}(A,B))$
relative to $S_5$ satisfies
\[
 \beta_{\mathbb Q}
   =2\,[\Pi(p_Bc^Aj)]\quad\hbox{in }
        H^1_{\rm cts}(G_{\mathbb Q},M).
 \tag{E109.8}
\]
Equations (E109.12)--(E109.14) retain the coordinate and inner
corrections proving this equality.  They show that every allowed
change of the affine origin or geometric paths preserves this
nonzero class.  This evaluates a necessary arithmetic term for
the mixed compact inputs of E.23.  Their cup obstructions and
central fibres remain on the active list; this statement concerns
the specified class-two actions and their coefficient maps.
The manuscript anchor is \S51.16, folios 733v--736v,
scans 172--178, through the compact factor of its generalized
Jacobian and the next approximation in Part E(c),(d).
\end{statement}
\begin{proof}
\emph{The exact central quotient and the coefficient left inverse.}
Filling the twelve punctures gives the continuous surjection of
geometric pro-five groups.  In a symplectic presentation of the
compact group its single relation is
$\prod_{r=1}^3[a_r,b_r]=1$.  Its class-two image is the central
relation $q=\sum_re_r\wedge f_r$.  The free affine group has
rank $17$; its class-two BCH coordinates and the quotient map are
therefore the actual groups $A\oplus B$ and $H\oplus Z$ with law
$(x,z)(y,w)=(x+y,z+w+x\wedge y/2)$.
One obtains the free class-two coordinates by the ordered Hall
normal form; its change to BCH coordinates has the inverse given
in E.105.  Quotienting by the central element $q$ proves the
compact assertion and all ranks in (E109.1).

The CM isogenies $\Phi,\Psi$ of (E103.13) have composites $[4]$.
On five-adic Tate modules they are inverse after division by
four.  Transfer the common CM operator $[i]$ on $E_+^3$ by these
maps.  Its square is $-1$.  Its two idempotents, formed by dividing
$1\pm i^{-1}[i]$ by two, give the stated integral direct sum.
The polarization pairs its two eigenspaces perfectly and is zero
on either one separately.  We may thus choose a symplectic CM
basis $(e_1,f_1,e_2,f_2,e_3,f_3)$, with all $e_r$ in $H^+$
and all $f_r$ in $H^-$.  The projector $\pi_M$ kills
$q\wedge H$.

For $D\in\operatorname{Hom}(H,Z)$ choose any $R$-linear lift
$\widetilde D:H\to\Lambda^2H$ and set
\[
 \operatorname{Alt}(\widetilde D)=
 \sum_{r=1}^3
  \bigl(e_r\wedge\widetilde D(f_r)
                -f_r\wedge\widetilde D(e_r)\bigr),
 \tag{E109.9}
\]
with the implicit inverse Tate twist provided by the pairing.
Changing the lift adds an element of $q\wedge H$.
The canonical tensor duality
$\operatorname{Hom}(H,Z)=H(-1)\otimes Z$ proves that the
result modulo $q\wedge H$ is independent of the symplectic
basis.  Consequently
\[
       \Pi(D)=\tfrac13\pi_M\operatorname{Alt}(\widetilde D)
 \tag{E109.10}
\]
is well-defined and equivariant.  Expanding a decomposable
$a\wedge b\wedge c$ in (E109.3) gives
$\operatorname{Alt}(\widetilde{f(\eta)})=3\eta$
modulo $q\wedge H$.  Expanding $\operatorname{ad}_h$ gives
$-2q\wedge h$ before any choice of quotient lift.  These two
identities prove $\Pi f=\pi_M$ and $\Pi\operatorname{ad}=0$,
including on continuous cochains.

For completeness the asserted injections are integral.  Let
$C_3:\Lambda^3H\to H(1)$ be symplectic contraction.
Then $C_3(q\wedge h)=2h$.  If $f(\eta)=0$, (E109.9) first
gives $3\eta\in q\wedge H$, hence $\eta=q\wedge h$.
Equation $f(q\wedge h)=-\operatorname{ad}_h$ then gives
$h\wedge x\in Rq$ for every $x$.  Reduction modulo five
forces $h=0$ modulo five: a nonzero decomposable two-vector has
rank two, whereas the alternating matrix of $q$ has rank six.
The same argument after removing the common power of five proves
$h=0$ over $R$.  The argument also proves that the inner map is
injective modulo five.  It therefore has a unit maximal minor,
so its image is a direct summand of the rank-$84$ Hom module.
This proves the freeness and rank in (E109.2), and the quotient
injection in (E109.3).  The certificate retains the full matrices
of these maps, with $\Pi f|_M=I_2$ and $\Pi\operatorname{ad}=0$
as literal matrix identities.

\emph{Filling the cusps with every coordinate correction retained.}
An automorphism of either BCH group is $R$-linear: its group
law and its commutators recover addition by removing the term
$[x,y]/2$, and its continuous integer powers recover scalar
multiplication.  Thus an arithmetic element acts as
\[
  g(x,z)=(\rho_A(g)x,\rho_B(g)z+D^A_gx),
  \qquad c^A_g=D^A_g\rho_A(g)^{-1},
 \tag{E109.11}
\]
and the corresponding compact formula defines $D^C_g,c^C_g$.
Multiplication of these block matrices proves
$c_{gh}=c_g+g c_h$ with the usual Hom-module action.  The actual
filling map in the chosen coordinates has the form
\[
       (x,z)\longmapsto(px,p_Bz+Jx)
       \quad(J\in\operatorname{Hom}_R(A,Z)).
 \tag{E109.12}
\]
We retain $J$.  Commuting this map with the actual arithmetic
actions, and then using the actual equivariant section $j$, gives
\[
 \begin{split}
  D^C_gp&=p_BD^A_g+J\rho_A(g)-\rho_Z(g)J,\\
  c^C_g&=p_Bc^A_gj+Jj-g(Jj).
 \end{split}
 \tag{E109.13}
\]
The second equality exhibits the exact continuous coboundary.
There is no inference from an associated graded action in this
step.

Changing an arithmetic section by a geometric cocycle $b_g$
changes its action to $\operatorname{Int}(b_g)g$.  In class two,
only the abelian component of $b_g$ contributes to this change:
\[
           c^A_g\longmapsto c^A_g+\operatorname{ad}_{b_g^{\rm ab}},
 \qquad c^C_g\longmapsto c^C_g+\operatorname{ad}_{p b_g^{\rm ab}}.
 \tag{E109.14}
\]
Changing BCH splittings adds a continuous Hom coboundary.
Equations (E109.10), (E109.13) and (E109.14) prove (E109.8)
and all the claimed transports.  Over $F$ the filled cusp section
is the ordinary point section at $b$; any discrepancy with the
filled $S_5$ is handled by (E109.14).

\emph{The normalization relating the action to the cycle.}
We use the integral Harris--Pulte comparison with its factor two:
Hain--Matsumoto, Theorem 3 and the proof in \S8, together with
their Magnus convention and Proposition 5.5
\cite{E109_HainMatsumoto}.  The Magnus map there is
$g\mapsto(x\mapsto g(\widetilde x)\widetilde x^{-1})$.
It agrees with (E109.11) on automorphisms acting trivially on $H$.
The resulting cycle class is twice this extended action class.
Here is also the comparison to the continuous group-ring
construction, making the coefficient map explicit.

For the augmentation ideal $I$ of the pro-five surface group,
the map
\[
        u\longmapsto \log_2(u)
           =(u-1)-\tfrac12(u-1)^2\pmod{I^3}
 \tag{E109.15}
\]
has the BCH product above; its central part sends a commutator
to the corresponding alternating tensor.  The surface relation
maps precisely to $q$.  A linear BCH splitting of $H$ therefore
gives a linear splitting of $I/I^3\to I/I^2$ whose extension
cocycle on the alternating central part is $c^C$.
The images of a lift and any other such lift differ by a linear
map, yielding exactly a coboundary.  The quotient by inner
conjugations is (E109.2).  Corollary 2.17 and Proposition 2.23
of \cite{E109_BisognoLiLittSrinivasan}, with their complete
comparison proof, give under these maps
\[
             2[c^C|_{G_F}]_{\rm out}
                =\bar f_*\nu(C),\qquad
  \nu(C)=\mu(C,b)\pmod{q\wedge H}.
 \tag{E109.16}
\]
All divisions in (E109.10) and (E109.15) are units of $R$.

\emph{The finite automorphisms project the entire class to $M$.}
The actual $F$-automorphisms
$(x,y,z)\mapsto(ix,y,z)$ and $(x,y,z)\mapsto(x,iy,z)$
generate a group $\Delta\simeq(\mathbb Z/4)^2$.
On the three maps $(f_0,f_1,f_2)$ of (E103.11), their two
CM exponents are respectively
\[
               (1,1),\qquad(2,1),\qquad(1,2).
 \tag{E109.17}
\]
For example $f_0(ix,y,z)=[i]f_0(x,y,z)$ and
$f_1(ix,y,z)=[-1]f_1(x,y,z)$; substituting the homogeneous
coordinates proves all six equalities.  The six characters of
$H$ are the three displayed characters and their negatives.
Adding any three distinct characters shows that the only
trivial characters in $\Lambda^3H$ are the sum of the three
displayed characters and its negative.  There are no trivial
characters in $H$.  It follows, with actual averaging projector
$16^{-1}\sum_{\delta\in\Delta}\delta$, that
\[
           (L/(q\wedge H))^\Delta=M.
 \tag{E109.18}
\]
The code checks all twenty exterior characters and this projector.

The class $\nu(C)$ is fixed by every automorphism of the curve.
This also follows directly on our action cocycles: transport by
a geometric automorphism conjugates their block matrices, which
changes the coefficient action by a Hom coboundary; a change
of based path adds the inner term (E109.14).  Thus averaging
the coefficient module averages the same cohomology class.
Equation (E109.18) and (E109.16) give
\[
        2[c^C|_{G_F}]_{\rm out}
              =\bar f_*\pi_{M*}\mu(C,b).
 \tag{E109.19}
\]
The module $M$ and the projector $\Pi$ are $G_{\mathbb Q}$-stable:
Galois either preserves or exchanges the two common CM
eigenspaces.  Define $\beta_{\mathbb Q}=2\Pi_*[c^C]_{\rm out}$.
The restriction of
$2[c^C]_{\rm out}-\bar f_*\beta_{\mathbb Q}$ is zero by
(E109.19) and $\Pi\bar f|_M=1$.  Applying corestriction to this
restriction gives multiplication by $[F:\mathbb Q]=4$ on the
original cohomology class.  The maps
are the finite coset sums on continuous cochains, so this
identity applies to the compact coefficients used here.
Multiplication by four is invertible.  This proves the global
equality in (E109.6), without a rational point on $C$ over
$\mathbb Q$.

\emph{Literal comparison with Bloch's quartic and its maps.}
Put $C_-:T_0^4+T_1^4=T_2^4$ and
$E_-:Y^2=X^3-X$.  The two $F$-isomorphisms used are
\[
 C\longrightarrow C_-,\quad[x:y:z]\longmapsto[x:y:sz],
 \qquad E_-\longrightarrow E_+,\quad(X,Y)\longmapsto(s^2X,s^3Y).
 \tag{E109.20}
\]
Write Bloch's maps on $z=1$, after substituting the first
isomorphism, as
\[
 \rho_1=(-y^2/s^2,x^2y/s^3),\qquad
 \rho_2=(-x^2/s^2,y^2x/s^3),\qquad
 \rho'=(x^2/y^2,s^2x/y^3).
 \tag{E109.21}
\]
These are the three maps of his \S3 to $E_-\times E_-\times E_+$.
Let $T=(0,0)$ on $E_+$; translation by $T$ is
$\tau_T(X,Y)=(1/X,-Y/X^2)$.  In the notation $\operatorname{sc}$
for the second map of (E109.20), direct substitution gives
\[
      \operatorname{sc}\rho_1=f_1,\qquad
      \operatorname{sc}\rho_2=-\tau_Tf_2,\qquad
      \rho'=[i]f_0.
 \tag{E109.22}
\]
Every equality extends from this dense chart to the proper
curves.  These equations prove that the induced Jacobian map
$\Phi_B:J_C\to E_+^3$ is $(f_{1*},-f_{2*},[i]f_{0*})$.
It is the actual E.103 isogeny followed by a permutation and
CM units, and its five-adic cohomology maps are isomorphisms.
They carry the pure CM projector in (E109.4) to Bloch's projector
of Hodge types $(3,0)+(0,3)$, with the same Tate twist.

We also retain the translation in (E109.22).  The three images
of $b$ are $(0,0),(-i,0),O$, all killed by two.  Hence, as
literal maps of curves,
\[
               [2]\rho=([2]\Phi_B)\circ i_b,
 \tag{E109.23}
\]
where the first two elliptic factors have been transported by
$\operatorname{sc}$.  Push this equality forward on the difference
of the curve and its negative.  On $H^3$ of an abelian threefold,
$[2]_*$ is multiplication by $8$: its pullback is multiplication
by $2^3$, and push--pull is multiplication by its degree $2^6$.
Functoriality of the Abel--Jacobi map therefore gives
$8\operatorname{AJ}(\rho(C)-[-1]\rho(C))=8\Phi_{B*}\mu(C,b)$.
Division by eight in $R$ proves the exact cycle-class comparison.
Thus no invariance under translation has been assumed.

\emph{A nonzero, nontorsion integral component.}
Bloch's original \S2, Proposition 2.1 and Corollary 2.2, and
\S3 project the finite correspondence test to the same pure CM
module.  In his \S4 the image $\beta_5$ of the cycle in
$H^1(G_{\mathbb Q(i)},M)$ is explicitly proved nonzero using
this test, and its nontorsion follows from the invariant test
at eleven \cite{E109_Bloch}.  We recompute the finite input:
\[
 \delta_p=\sum_{\substack{u\in\mathbb F_{p^2}\\u+u^p=1}}
 u^{(p^2-1)/2}
       \bigl(u^{(p-1)/2},u^{(p-3)/4}\bigr)
             \quad\text{in }E'(\mathbb F_p),\quad p\equiv3\pmod4.
 \tag{E109.24}
\]
The sign multiplies the point using its group law.  In the
correspondence of Bloch's Lemma 3.2 the detecting point is
$4\delta_p$; (E109.7) thus supplies its nonzero five-primary
part.  Our code implements all field operations in
$\mathbb F_p[i]/(i^2+1)$, verifies every point and addition,
and checks exact orders by prime divisors.  It also independently
checks formula 3.6 of the source, retaining its two-torsion
difference from (E109.24), for every $p\equiv3\pmod4$ up to $103$.

Here is the torsion argument at the actual coefficient group.
For the prime over eleven of $\mathbb Q(i)$, Frobenius on
$M\otimes\mathbb Q_5/R$ is the scalar $-11$.
Its difference from one is the five-adic unit $-12$.
Thus this module has no $G_{\mathbb Q(i)}$-invariants.
The exact cohomology sequence for
$0\to M\to M\otimes\mathbb Q_5\to M\otimes\mathbb Q_5/R\to0$
makes the map from $H^1(G_{\mathbb Q(i)},M)$ into rational
cohomology injective.  The nonzero class $\beta_5$ is therefore
nontorsion.  Restriction to $F$ preserves this property: the
restriction--corestriction composite is multiplication by two,
which is invertible.  Equations (E109.20)--(E109.23) identify
its restriction with $\pi_{M*}\mu(C,b)$ through an actual
five-adic isomorphism.  Equations (E109.19) and (E109.10)
now prove every assertion of (E109.6), including nonvanishing
after rational scalar extension.

The integral coefficient map with left inverse is essential in
this last inference: it gives a left inverse on continuous
cohomology itself.  The finite-field and matrix certificates
support the displayed calculations; the cycle comparison,
cohomological transports and original theorem applications
above supply the proof for the actual arithmetic action.
\end{proof}
